\section{Related Work}
\label{sec:related}

\para{Lies vs.\ honest errors.}
\mask~\citep{ren2025mask} defines a lie as a statement that contradicts the model's belief, where belief is elicited by consistent neutral answers, and it separates honesty from accuracy behaviourally. It reports a lie rate of 23.5\% for \llama under pressure, but it does not test probes. Liars' Bench~\citep{kretschmar2025liars} collects 72,863 lies and honest answers from four models and verifies beliefs with a noise-floor rule. \citet{cooney2026didyoulie} build a \triviaqa testbed in which lie-inducing system prompts are applied only to questions the model answers correctly. Both evaluate detectors against honest answers that are correct. WACK~\citep{simhi2024distinguishing} splits hallucinations into errors without knowledge (HK$-$) and errors despite knowledge (HK$+$) and uses a sampling-based knowledge test similar to ours. \citet{liu2023cognitive} find that most disagreements between probes and outputs reflect calibration rather than deception. Building on these definitions, we cross knowledge with incentive on the same questions and measure the share of each cause.

\para{White-box lie detectors.}
Representation engineering~\citep{zou2023representation} introduced honest/dishonest persona contrasts, which \citet{goldowsky2025detecting} turned into the standard deception probe: logistic regression on mean response-token activations. They note that their probe also fires on honest answers in deception-salient contexts. Truth probes trained on true/false statements~\citep{azaria2023internal,marks2023geometry,burger2024truth} are a second family; \citet{thormann2026limits} argues that they are falsity detectors. \citet{natarajan2026oneprobe} show that the choice of training prompt explains most of the variance in probe performance, and \citet{moustafa2026beyond} show that the lie typology of the training data dominates detectability. \citet{luikham2026asymmetries} find only partial transfer between instructed and spontaneous deception probes, and \citet{huan2025canllmslie} report that lies form a separate cluster from hallucinations in \llama, using prompts that contain the word ``lie''. Black-box alternatives ask follow-up questions~\citep{pacchiardi2023catch}. Unlike these studies, we test lie detectors against false-but-ignorant negatives and against probes that see only the prompt.

\para{Hallucination detectors.}
Hallucination detection uses token probabilities, self-evaluation such as P(True) and P(IK)~\citep{kadavath2022language}, semantic entropy over samples~\citep{kuhn2023semantic} and its probe-based approximation~\citep{kossen2024semantic}, and correctness probes at the answer token~\citep{orgad2025llms}. \citet{cheang2025internal} show that internal states mostly reflect whether knowledge recall happens, not whether the output is true. This predicts that a lie, which involves full recall of the right answer, will look like a correct answer to such detectors. We test this prediction directly.
